\section*{Capítulo \ref{ch:b1:observing-the-sky} --- Telescopios, detectores y fotometría}

\begin{solution}{exo:b1:observing-the-sky:1}
Una diferencia de 5 magnitudes equivale a un factor de 100: la estrella de magnitud 1.0
envía 100 veces más \omterm{def:b1:observing-the-sky:flux}{flujo}. Rigel y Polaris: $10^{0.4\times(2.02 - 0.13)} =
10^{0.756} = 5.70$; Rigel, con la magnitud menor, es la más brillante.
\end{solution}

\begin{solution}{exo:b1:observing-the-sky:2}
$d = 1/0.19495 = \qty{5.13}{\parsec}$, es decir, $5.1295\times206265 =
\qty{1.058e6}{\au}$ y $5.1295\times3.2616 = \qty{16.7}{\lightyear}$.
\end{solution}

\begin{solution}{exo:b1:observing-the-sky:3}
$1.22\times\num{550e-9}/2.4 = \qty{2.80e-7}{rad} = \ang{;;0.058}$, unas siete veces más
fino que un \omterm{def:b1:observing-the-sky:seeing}{seeing} de \ang{;;0.40}: en el espacio el telescopio alcanza su límite de
difracción. Para \qty{6.5}{m} a \qty{2.0}{\micro\metre}:
$1.22\times\num{2.0e-6}/6.5 = \qty{3.75e-7}{rad} = \ang{;;0.077}$.
\end{solution}

\begin{solution}{exo:b1:observing-the-sky:4}
$B - V$ sube de Rigel ($-0.03$) a Altair (0.22), Polaris (0.60) y Betelgeuse (1.85): la
más caliente es Rigel, luego Altair, Polaris y Betelgeuse. La distancia atenúa ambas bandas
por el mismo factor, así que $B - V$ no cambia por ella (mientras no haya polvo que enrojezca
la luz en el camino).
\end{solution}

\begin{solution}{exo:b1:observing-the-sky:5}
$X_1 = \sec\ang{30} = 1.1547$, $X_2 = \sec\ang{60} = 2.0000$.
$k = (6.590 - 6.412)/(2.0000 - 1.1547) = 0.178/0.8453 = \num{0.211}$ por \omterm{def:b1:observing-the-sky:extinction}{masa de aire}, y
$m_0 = 6.412 - 0.2106\times1.1547 = 6.169$.
\end{solution}

\begin{solution}{exo:b1:observing-the-sky:6}
$\mathrm{SNR} = \sqrt{400} = 20$. La \omterm{def:b1:observing-the-sky:detector}{relación señal-ruido} crece como $\sqrt t$: para 100,
$t = \qty{10}{s}\times(100/20)^2 = \qty{250}{s}$. Un error de \num{0.01} en la magnitud
exige $\mathrm{SNR} = 1.086/0.01 = 108.6$, así que
$t = 10\times(108.6/20)^2 = \qty{295}{s}$.
\end{solution}

\begin{solution}{exo:b1:observing-the-sky:7}
La distancia es $d = 1/0.00378 = \qty{264.6}{\parsec}$, y la \omterm{def:b1:observing-the-sky:absolute}{magnitud absoluta}
$M = 0.13 - 5\log_{10}26.46 = 0.13 - 7.11 = -6.98$. $L/L_\odot = 10^{0.4\times(4.83 + 6.98)} =
10^{4.724} = \num{5.3e4}$ en V.
\end{solution}

\begin{solution}{exo:b1:observing-the-sky:8}
$\pi\times1.2^2 = \qty{4.52}{m^2}$, y $25/4.52 = 5.5$: el telescopio Webb recoge el mismo
número de fotones en una exposición 5.5 veces más corta. Un disco completo de
\qty{6.5}{m} de diámetro tendría $\pi\times3.25^2 = \qty{33.2}{m^2}$; el espejo real tiene
$25/33.2 = \qty{75}{\percent}$ de él, porque 18 hexágonos dejan sin cubrir el borde de ese
disco, y los huecos entre segmentos y la abertura central no recogen nada.
\end{solution}

\begin{solution}{exo:b1:observing-the-sky:9}
(a) $\qty{10}{\parsec} = \qty{3.0857e17}{m}$, y $F = \num{3.828e26}/(4\pi
\times(\num{3.0857e17})^2) = \num{3.828e26}/\num{1.1965e36} =
\qty{3.20e-10}{W.m^{-2}}$, $1361/\num{3.20e-10} = \num{4.3e12}$ veces menos que en la Tierra.
(b) $5\log_{10}(d/\qty{10}{\parsec}) = -26.74 - 4.83 = -31.57$, así que
$d = \qty{10}{\parsec}\times10^{-6.314} = \qty{4.853e-6}{\parsec}$, es decir,
$\num{4.853e-6}\times206265 = \qty{1.001}{\au}$.
\end{solution}

\begin{solution}{exo:b1:observing-the-sky:10}
Por encima de la atmósfera el \omterm{def:b1:observing-the-sky:flux}{flujo} es $F_0$; por debajo, $F = F_0e^{-\tau_0X}$. Entonces
$m - m_0 = -2.5\log_{10}(F/F_0) = 2.5\log_{10}e\;\tau_0X$, que es $kX$ con
$k = 2.5\log_{10}e\;\tau_0 = 1.0857\,\tau_0$. Para $k = 0.160$:
$\tau_0 = 0.160/1.0857 = 0.147$. En el \omterm{def:b1:spherical-astronomy:celestial-sphere}{cenit} llega $e^{-0.147} = 0.863$ de la luz, el
\qty{86}{\percent}; a $X = 2.31$, $e^{-0.147\times2.31} =
e^{-0.340} = 0.712$, el \qty{71}{\percent}. A $z = \ang{85}$, $\sec z =
11.5$; pero el modelo plano-paralelo hace infinito el camino en el horizonte, mientras que el
camino a través de una atmósfera curva sigue siendo finito, y además la refracción curva el
rayo: cerca del horizonte la fórmula sobrestima la \omterm{def:b1:observing-the-sky:extinction}{masa de aire} y sus errores crecen deprisa.
\end{solution}

\begin{solution}{exo:b1:observing-the-sky:11}
(a) $\mathrm{SNR} = Rt/\sqrt{(R + B)t}$, así que $t = \mathrm{SNR}^2(R + B)/R^2
= 100\times2050/2500 = \qty{82}{s}$. (b) Reducir a la mitad el \omterm{def:b1:observing-the-sky:seeing}{seeing} divide el área de
la imagen, y con ella $B$, por 4: $B = 500$ y $t = 100\times550/2500 =
\qty{22}{s}$. (c) Para $B \gg R$, $t \approx \mathrm{SNR}^2B/R^2$. La tasa de la fuente $R$
crece como $D^2$; la tasa del cielo $B$ crece como el área colectora $D^2$ por el área del
cielo en la imagen, $\theta^2$. Por tanto, $t \propto
\theta^2D^2/D^4 = \theta^2/D^2$: reducir a la mitad el \omterm{def:b1:observing-the-sky:seeing}{seeing} vale tanto como duplicar el
diámetro.
\end{solution}

\begin{solution}{exo:b1:observing-the-sky:12}
(a) $N = 1200\times100 = \num{120000}$ líneas; en el primer orden el poder resolutor es
\num{1.2e5}, y $\Delta\lambda = 550/(\num{1.2e5}) =
\qty{4.6e-3}{nm}$. (b) $v = c\,\Delta\lambda/\lambda = c/(\num{1.2e5}) =
\qty{2.5}{km.s^{-1}}$. (c) $\Delta\lambda = 550\times29.78/\num{299792} =
\qty{0.055}{nm}$, unos 12 elementos de resolución. El observador debe calcular, para la fecha
y la hora de la observación, la componente de la velocidad orbital de la Tierra hacia la
estrella, y restarla, para obtener la velocidad de la estrella respecto del Sol.
\end{solution}

\begin{solution}{pb:b1:observing-the-sky:1}
\textbf{1.} $d = 1/0.37921 = \qty{2.637}{\parsec}$.

\textbf{2.} $2.637\times206265 = \qty{5.44e5}{\au}$;
$2.637\times\num{3.0857e16} = \qty{8.14e16}{m}$.

\textbf{3.} $2.637\times3.2616 = \qty{8.60}{\lightyear}$: la luz salió de Sirio hace 8.6 años.

\textbf{4.} La dirección oscila $2p = \ang{;;0.758}$, unas dos veces el \omterm{def:b1:observing-the-sky:seeing}{seeing}. La oscilación
se mide a partir de la posición del centro de la imagen, no de su borde: promediado sobre
muchos fotones, el centro se halla con una fracción pequeña de la anchura de la imagen, tanto
más precisamente cuantos más fotones se cuenten.

\textbf{5.} $\Delta d/d = \Delta p/p = 0.001/0.37921 = \qty{0.26}{\percent}$,
es decir, $8.60\times0.0026 = \qty{0.023}{\lightyear}$.

\textbf{6.} $10^{0.4\times(0.03 + 1.46)} = 10^{0.596} = 3.94$.

\textbf{7.} $10^{0.4\times(26.74 - 1.46)} = 10^{10.112} = \num{1.29e10}$.

\textbf{8.} $10^{0.4\times(8.44 + 1.46)} = 10^{3.96} = \num{9120}$.

\textbf{9.} $z = \ang{90} - \ang{82.09} = \ang{7.91}$, $X = \sec\ang{7.91} =
1.0096$; extinción $0.15\times1.0096 = 0.151$ magnitudes; medida
$m = -1.46 + 0.151 = -1.31$.

\textbf{10.} $z = \ang{60}$, $X = 2.000$; extinción de 0.30 magnitudes,
$m = -1.16$. La fracción transmitida es $10^{-0.4\times0.30} = 0.759$:
\qty{24}{\percent} de la luz se pierde.

\textbf{11.} Al salir $z = \ang{90}$ y $\sec z$ es infinito, mientras que el camino real
a través de la atmósfera curva es largo pero finito; la refracción además curva el rayo. Cerca
del horizonte la extinción es grande y cambia deprisa, y la corrección plano-paralela es
incorrecta.

\textbf{12.} $\num{3.0e4}\times10^{0.4\times(10.0 - 8.44)} = \num{3.0e4}
\times10^{0.624} = \num{3.0e4}\times4.21 = \num{1.26e5}$ fotoelectrones por segundo.

\textbf{13.} Una \omterm{def:b1:observing-the-sky:detector}{relación señal-ruido} de 1000 exige $N = 1000^2 = \num{1.0e6}$
fotoelectrones: $t = \num{1.0e6}/\num{1.26e5} = \qty{7.9}{s}$. El error en la magnitud es
$1.086/1000 = \num{0.0011}$.

\textbf{14.} $\num{3.0e4}\times10^{0.4\times(10.0 + 1.46)} = \num{3.0e4}
\times10^{4.584} = \num{1.15e9}$ fotoelectrones por segundo.

\textbf{15.} $0.20\times\num{1.15e9} = \num{2.30e8}$ por segundo en el píxel más brillante,
que se llena en $\num{1.0e5}/\num{2.30e8} =
\qty{4.3e-4}{s}$, menos de medio milisegundo.

\textbf{16.} Las alas de Sirio dan $\num{1.15e9}\times10^{-4} =
\num{1.15e5}$ por segundo; Sirio~B da $0.20\times\num{1.26e5} =
\num{2.52e4}$ por segundo. El resplandor de Sirio es 4.6 veces más brillante que el pico de
Sirio~B.

\textbf{17.} El píxel recibe $\num{1.15e5} + \num{2.52e4} = \num{1.40e5}$
por segundo y se llena en $\num{1.0e5}/\num{1.40e5} = \qty{0.71}{s}$. En
\qty{0.5}{s}: señal $\num{2.52e4}\times0.5 = \num{1.26e4}$, total
$\num{7.01e4}$, ruido $\sqrt{\num{7.01e4}} = 265$, así que $\mathrm{SNR} =
\num{1.26e4}/265 = 48$. Sirio~B es detectable en lo que respecta al ruido de fotones; lo que
lo hace difícil es que el resplandor debe conocerse y restarse con una precisión muy superior a
una quinta parte de su valor, mientras que el núcleo de Sirio se satura en menos de un milisegundo.

\textbf{18.} $m - M = 5\log_{10}(2.637/10) = 5\times(-0.579) = -2.89$.

\textbf{19.} $M = -1.46 + 2.89 = 1.43$.

\textbf{20.} $d = 1/0.13023 = \qty{7.679}{\parsec}$;
$m - M = 5\log_{10}0.7679 = -0.57$, así que $M = 0.03 + 0.57 = 0.60$.

\textbf{21.} Vega, con la \omterm{def:b1:observing-the-sky:absolute}{magnitud absoluta} menor, por el factor
$10^{0.4\times(1.43 - 0.60)} = 10^{0.332} = 2.15$. Sirio parece más brillante solo porque
está más cerca.

\textbf{22.} $m = M_\odot + (m - M) = 4.83 - 2.89 = 1.94$: el Sol sería una estrella bastante
brillante a simple vista.

\textbf{23.} $5\log_{10}(d/\qty{10}{\parsec}) = 6.0 - 1.43 = 4.57$, así que
$d = 10\times10^{0.914} = \qty{82.0}{\parsec}$, unos \qty{267}{\lightyear}.

\textbf{24.} $L/L_\odot = 10^{0.4\times(4.83 - 1.43)} = 10^{1.36} = 22.9$:
en la banda V, Sirio supera en brillo a \textbf{unos 23 soles}.
\end{solution}
